%% ============================================================
%%  Chapter 4 — Structured Output, Validation & Self-Healing
%%  Production-Grade Agentic AI: LangGraph + Python Frameworks
%% ============================================================

\chapterquote{If the contract between components is vague, the
system will fail in vague ways.}{Mary Poppendieck, adapted}

By the end of Chapter~1 you had a four-node loop running in production: a
Perceive node that redacts PII, a Reason node that calls an LLM via LiteLLM and
Instructor, an Action node that dispatches tools, and a Reflect node that
chooses among five enum decisions. Chapter~2 wrapped that loop in termination
policy and failure guards. This chapter tightens the remaining gap: what happens
between the raw text the model emits and the strongly typed decisions the rest
of the runtime expects.

Unstructured model output is the single biggest source of low-grade, expensive
instability in agent systems. ``Low-grade'' because runs usually complete; the
system only fails on edge cases. Expensive because each failure burns full model
calls, tool calls, and engineer time to debug. The correct response is not
``better prompts'' but a runtime-enforced contract: every node that consumes
model output must receive a Pydantic-validated object, and every parse failure
must either repair itself within a bounded number of retries or escalate in a
structured, observable way.

This chapter implements that contract. You will define strict Pydantic~v2 models
for the three Reason outputs, integrate Instructor's \code{patch()} modes as a
first-class configuration, introduce a dedicated \code{ValidationNode} into the
graph, and wire a repair loop that injects validation errors back into the next
model call as beliefs. When repair fails, a rejection branch marks the run as a
validation failure and surfaces the raw output and error to Langfuse for
analysis.

%% ------------------------------------------------------------
\section{Pydantic Output Models as Contracts}
\label{sec:output-models}
%% ------------------------------------------------------------

Chapter~1 sketched the discriminated union of model intents the Reason node can
produce: a tool call request, a final answer, or a reasoning trace. In this
chapter you make that union the explicit contract between the model call and the
rest of the runtime.

Each of the three variants gets its own Pydantic~v2 \code{BaseModel}. There is
no single ``assistant response'' model with optional fields for every possible
shape. Optional fields create mixed states that are syntactically valid and
semantically meaningless: a record with both a tool name and a final answer, or
neither. A discriminated union, keyed by a \code{type} field, rules those states
out.

Listing~\ref{lst:output-models} shows the final models. Two details matter for
production: strict mode and confidence bounds. Strict mode prevents Pydantic
from silently coercing values (\code{"1"} into \code{1}), which is attractive in
form handling but dangerous when parsing model output. Confidence bounds force
any model calibration bugs into validation failures where they belong.

\begin{lstlisting}[language=Python,
  caption={Strict Pydantic models for the three Reason output variants},
  label={lst:output-models}]
# src/agent/output_models.py

class ToolRequest(BaseModel):
    model_config = {"strict": True}
    type:      Literal["tool_request"] = "tool_request"
    tool_name: str
    arguments: dict[str, Any]

class FinalAnswer(BaseModel):
    model_config = {"strict": True}
    type:       Literal["final_answer"] = "final_answer"
    content:    str
    confidence: float

    @field_validator("confidence")
    @classmethod
    def confidence_in_range(cls, v: float) -> float:
        if not 0.0 <= v <= 1.0:
            raise ValueError(
                f"confidence must be between 0.0 and 1.0, got {v}")
        return v

class ReasoningTrace(BaseModel):
    model_config = {"strict": True}
    type:  Literal["reasoning_trace"] = "reasoning_trace"
    trace: str

ReasonOutput = Annotated[
    ToolRequest | FinalAnswer | ReasoningTrace,
    Field(discriminator="type"),
]
\end{lstlisting}

These models are imported by the Reason and Validation nodes, not redefined. The
union alias \code{ReasonOutput} is the only response model used when calling the
LLM via Instructor. That keeps the number of model schemas in your system small
and explicit.

\begin{warningbox}
Do not introduce a catch-all \code{"other"} variant into the union. If the model
produces a shape you did not anticipate, that is a validation failure. It should
trigger the repair loop or escalation, not silently flow downstream as an
untyped blob.
\end{warningbox}

%% ------------------------------------------------------------
\section{Choosing an Instructor Mode}
\label{sec:instructor-modes}
%% ------------------------------------------------------------

Instructor wraps a LiteLLM client so that calls to \code{.create()} accept a
\code{response\_model=} argument. Under the hood, it can use three strategies to
coerce the model into structured output: native tool calling, native JSON mode,
or JSON wrapped in a markdown code block. The mode you choose is a property of
both the model provider and the reliability you need.

The three modes are captured as an enum in Listing~\ref{lst:instructor-config}.
\code{TOOLS} uses the provider's function-calling API when available. \code{JSON}
uses explicit JSON mode when the provider supports it. \code{MD\_JSON} is a
last resort for models that insist on wrapping JSON in a code fence.

\begin{lstlisting}[language=Python,
  caption={Selecting an Instructor mode per model family},
  label={lst:instructor-config}]
# src/agent/instructor_config.py

class InstructorMode(str, Enum):
    TOOLS   = "tools"
    JSON    = "json"
    MD_JSON = "md_json"

_MODE_MAP: dict[InstructorMode, instructor.Mode] = {
    InstructorMode.TOOLS:   instructor.Mode.TOOLS,
    InstructorMode.JSON:    instructor.Mode.JSON,
    InstructorMode.MD_JSON: instructor.Mode.MD_JSON,
}

MODEL_TO_INSTRUCTOR_MODE: dict[str, InstructorMode] = {
    "gpt-4o":                 InstructorMode.TOOLS,
    "gpt-4o-mini":            InstructorMode.TOOLS,
    "claude-3-5-sonnet":      InstructorMode.TOOLS,
    "mistral/mistral-large":  InstructorMode.JSON,
    "ollama/llama3":          InstructorMode.MD_JSON,
}


def make_instructor_client(
    model: str,
    mode: InstructorMode | None = None,
) -> instructor.Instructor:
    """Patch a LiteLLM client with the right Instructor mode."""
    base_client = litellm.Client()
    resolved = mode or MODEL_TO_INSTRUCTOR_MODE.get(
        model, InstructorMode.JSON
    )
    return instructor.from_litellm(
        client=base_client, mode=_MODE_MAP[resolved]
    )
\end{lstlisting}

In prose, the decision flow is simple:
\begin{enumerate}
  \item If the provider supports function calling, use \code{TOOLS}.
  \item Else if it supports JSON mode, use \code{JSON}.
  \item Otherwise, fall back to \code{MD\_JSON} and treat parse failures as
        expected until you can switch models.
\end{enumerate}

Chapter~9 returns to this enum when it replaces the hand-written token estimate
with LiteLLM's actual usage callbacks. The same mapping is used to compute cost
per mode.

%% ------------------------------------------------------------
\section{Validation as a Graph Node}
\label{sec:validation-node}
%% ------------------------------------------------------------

In Chapter~1 the Reason node called Instructor directly and wrote the parsed
object into \code{last\_action}. Validation lived inline inside the node. That is
convenient but wrong for three reasons:

\begin{itemize}
  \item you cannot route on validation failures without hiding control flow
        inside the node function,
  \item you cannot observe validation separately from reasoning, and
  \item you cannot reuse the same validation logic for other nodes.
\end{itemize}

The fix is to introduce a dedicated \code{ValidationNode} into the graph. The
Reason node is responsible only for calling the model and capturing the raw
output string. The Validation node is responsible for parsing that string into a
\code{ReasonOutput}, recording failures, and choosing between the happy path and
repair.

Three private fields are added to the state for this purpose:

\begin{itemize}
  \item \code{\_raw\_model\_output: str | None} --- the last raw JSON string the
        model produced,
  \item \code{\_validation\_error: str | None} --- the last Pydantic
        \code{ValidationError} string, and
  \item \code{\_validation\_attempt: int} --- the number of consecutive failed
        validation attempts for the current cycle.
\end{itemize}

Listing~\ref{lst:validation-node} shows the node. It is a pure
\code{(state) -> dict} callable like the other nodes, decorated with
\code{@observe} so Langfuse captures it as a separate span. On success it writes
\code{last\_action} as a fully parsed union and clears the error fields. On
failure it increments \code{\_validation\_attempt}, records the error, and leaves
\code{last\_action} untouched.

\begin{lstlisting}[language=Python,
  caption={Validation as its own node between Reason and Action},
  label={lst:validation-node}]
# src/agent/nodes/validation.py

@observe(name="validation_node")
def validation_node(state: AgentState) -> dict:
    raw = state.get("_raw_model_output")
    if raw is None:
        # Nothing to validate; this is a wiring error in the graph.
        return {}

    attempt = int(state.get("_validation_attempt", 0))
    try:
        parsed = ReasonOutput.model_validate_json(raw)
    except ValidationError as exc:
        error_str = str(exc)
        _emit_validation_failure_event(
            error=error_str,
            node_name="validation_node",
            attempt_number=attempt,
            model=state.get("_config_model", "unknown"),
            output_type="unknown",
            raw_output=raw,
        )
        return {
            "_validation_error":   error_str,
            "_validation_attempt": attempt + 1,
        }

    # Success: clear error fields and update last_action
    return {
        "last_action":         parsed.model_dump(),
        "_validation_error":   None,
        "_validation_attempt": 0,
    }
\end{lstlisting}

The output type is captured as \code{parsed.type} later in the routing logic for
reporting. At this point you only know that validation succeeded.

%% ------------------------------------------------------------
\section{A Bounded Repair Loop}
\label{sec:repair-loop}
%% ------------------------------------------------------------

A single validation failure is not a reason to give up. Models occasionally
hallucinate fields, omit required keys, or emit extra commentary before or after
JSON. Instructor's built-in retry loop catches many of these cases, but you
still need to handle the worst cases in the runtime. The pattern is simple:
parse fails, error is injected as a belief, Reason is called again, and validation
retries---up to a maximum.

The maximum must not be a magic number buried inside the node. It is part of the
termination policy. Chapter~2 introduced \code{TerminationPolicy} with fields
for \code{max\_cycles}, \code{token\_ceiling}, \code{wall\_clock\_timeout}, and
\code{cost\_cap}. Here you add a new field \code{validation\_max\_retries} with a
sane default and a validation rule.

\begin{lstlisting}[language=Python,
  caption={Extending the termination policy with a validation retry ceiling},
  label={lst:termination-policy-validation}]
# src/agent/termination.py

class TerminationPolicy(BaseModel):
    max_cycles:            int
    token_ceiling:         int
    wall_clock_timeout_s:  int
    cost_cap_usd:          float
    validation_max_retries: int = 2

    @field_validator("validation_max_retries")
    @classmethod
    def retries_must_be_positive(cls, v: int) -> int:
        if v < 1:
            raise ValueError(
                "validation_max_retries must be >= 1, got {v}")
        return v
\end{lstlisting}

The routing logic after the Validation node reads this policy from state (where
it was written at graph entry) and decides between pass, retry, and escalate.
Retrying consists of two steps: incrementing the attempt counter and appending a
repair belief summarising the error.

\begin{lstlisting}[language=Python,
  caption={Building a repair belief and routing based on the retry ceiling},
  label={lst:repair-loop}]
# src/agent/nodes/validation.py


def _build_repair_belief(error: str, attempt: int) -> str:
    return (
        f"[validation-retry attempt={attempt}] "
        f"ValidationError: {error}"
    )


def route_after_validation(state: AgentState) -> str:
    policy = TerminationPolicy.model_validate(
        state.get("_policy", {})
    )
    error   = state.get("_validation_error")
    attempt = int(state.get("_validation_attempt", 0))

    if error is None:
        return "pass"
    if attempt < policy.validation_max_retries:
        return "retry"
    return "escalate"
\end{lstlisting}

The belief injection itself happens when you handle the \code{"retry"} branch in
the graph wiring. The next Reason call will see the most recent validation error
in the \code{beliefs} list and can adjust its behaviour accordingly.

\begin{lstlisting}[language=Python,
  caption={Injecting the validation error into beliefs before retrying},
  label={lst:repair-belief-injection}]
# src/agent/graph.py (excerpt)


def _inject_validation_belief(state: AgentState) -> dict:
    error   = state.get("_validation_error")
    attempt = int(state.get("_validation_attempt", 0))
    if not error:
        return {}
    return {
        "beliefs": [
            _build_repair_belief(error=error, attempt=attempt)
        ]
    }
\end{lstlisting}

Chapter~2's token budget logic already accounts for each Reason call. You do not
need any special-case cost accounting for retries: they are just extra cycles.

%% ------------------------------------------------------------
\section{Escalation on Ceiling Breach}
\label{sec:validation-escalation}
%% ------------------------------------------------------------

When validation has failed \code{validation\_max\_retries} times in a row for the
same raw output, further retries are wishful thinking. At that point the runtime
needs to surface the failure explicitly and stop.

Chapter~1 introduced the Reflect node and its five enum decisions:
\code{CONTINUE}, \code{REVISE}, \code{RETRY}, \code{TERMINATE}, and
\code{ESCALATE}. Escalation here reuses that same decision. The Validation node
does not talk to the Reflect node directly; instead it sets state in a way that
causes the existing routing logic to choose the escalation branch.

Listing~\ref{lst:validation-escalation} shows a small helper for the escalation
update. It writes a structured summary into \code{final\_answer}, sets
\code{reflect\_decision} to \code{"ESCALATE"}, and leaves the error fields
intact for diagnostics.

\begin{lstlisting}[language=Python,
  caption={Marking a run as escalated due to validation failure},
  label={lst:validation-escalation}]
# src/agent/nodes/validation.py

@dataclass
class ValidationEscalationRecord:
    node_name:        str
    attempt_count:    int
    last_error:       str
    raw_model_output: str | None
    model:            str
    timestamp_utc:    str = field(
        default_factory=lambda:
            datetime.now(timezone.utc).isoformat()
    )


def _build_escalation_update(
    state: AgentState,
    error: str,
    attempt: int,
    model: str,
) -> dict:
    record = ValidationEscalationRecord(
        node_name="validation_node",
        attempt_count=attempt,
        last_error=error,
        raw_model_output=state.get("_raw_model_output"),
        model=model,
    )
    summary = (
        f"Validation failed after {attempt} retries. "
        f"Last error: {error[:200]}"
    )
    return {
        "reflect_decision": "ESCALATE",
        "final_answer":     summary,
        "_validation_error":   error,
        "_validation_attempt": attempt,
        "beliefs": [
            f"[validation-escalated] {summary}",
        ],
        "cycle_traces": [
            {"validation_escalation": asdict(record)},
        ],
    }
\end{lstlisting}

The Validation node calls this helper in its failure branch when
\code{route\_after\_validation} has returned \code{"escalate"}. The Reflect node
then sees \code{reflect\_decision == "ESCALATE"} and Chapter~1's
\code{route\_from\_reflect} function routes directly to \code{END}. The
\code{termination\_status} in the \code{CycleTrace} is set to \code{"FAILURE"} by
Chapter~2's helpers, distinguishing this path from voluntary human escalation.

%% ------------------------------------------------------------
\section{Langfuse Events for Validation Failures}
\label{sec:validation-langfuse}
%% ------------------------------------------------------------

Langfuse spans around Reason and Validation already tell you that validation
failed, how long it took, and how many retries were attempted. To understand
\emph{why} validation failed across a fleet of agents and models you need more
structure: an event with a consistent schema that you can aggregate by model,
output type, and error class.

This section adds a small helper that emits a Langfuse event every time
validation fails. Spans capture timing; events capture domain-specific facts.

\begin{lstlisting}[language=Python,
  caption={Emitting a structured Langfuse event on validation failure},
  label={lst:langfuse-event}]
# src/agent/nodes/validation.py


def _emit_validation_failure_event(
    error: str,
    node_name: str,
    attempt_number: int,
    model: str,
    output_type: str,
    raw_output: str | None,
) -> None:
    langfuse_context.update_current_observation(
        metadata={
            "event":            "validation_failure",
            "validation_error": error,
            "node_name":        node_name,
            "attempt_number":   attempt_number,
            "model":            model,
            "output_type":      output_type,
            "raw_output_len":   len(raw_output or ""),
        }
    )
\end{lstlisting}

The helper is called in the Validation node's \code{except ValidationError}
branch. For simple deployments this \code{metadata} record is enough: Langfuse
Studio can filter observations where \code{event == "validation\_failure"} and
show you error distributions by model and output type.

For larger fleets you can extend this into a dedicated Langfuse event type and a
data-warehouse export. The schema does not change; only the plumbing does.

\begin{notebox}
As in Chapter~1's tests, set \code{LANGFUSE\_ENABLED=false} in your test
configuration so that unit tests for the Validation node do not make network
calls. The helper function becomes a no-op when the SDK is disabled.
\end{notebox}

%% ------------------------------------------------------------
\section{Graph Wiring and Tests}
\label{sec:validation-graph}
%% ------------------------------------------------------------

The last step is to insert the Validation node into the graph and exercise the
repair and escalation paths in tests.

On the graph side, Reason writes \code{\_raw\_model\_output} instead of a parsed
\code{last\_action}. Validation consumes that field and writes \code{last\_action}
on success. The edge from Reason now points to Validation; the edge from
Validation points to Action or Reason depending on the routing function.

\begin{lstlisting}[language=Python,
  caption={Updated StateGraph wiring with a dedicated Validation node},
  label={lst:graph-ch4}]
# src/agent/graph.py (excerpt)

b.add_node("validation_node", validation_node)

b.add_edge("perceive_node", "reason_node")
b.add_edge("reason_node", "validation_node")

b.add_conditional_edges(
    "validation_node",
    route_after_validation,
    {
        "pass":     "action_node",
        "retry":    "reason_node",
        "escalate": END,
    },
)
\end{lstlisting}

Chapter~28 will add a full integration test suite for the graph. For now you
focus on unit tests that treat Validation as a pure function of state. Each test
constructs a minimal \code{AgentState}, populates \code{\_raw\_model\_output}, and
asserts on the node's return value.

\begin{lstlisting}[language=Python,
  caption={Node-level tests for validation, repair, and escalation},
  label={lst:validation-tests}]
# tests/test_validation.py

os.environ.setdefault("LANGFUSE_ENABLED", "false")


def test_valid_output_passes() -> None:
    state = make_test_state(
        _raw_model_output=FinalAnswer(
            content="done", confidence=0.9
        ).model_dump_json()
    )
    result = validation_node(state)
    assert result["last_action"]["type"] == "final_answer"
    assert result["_validation_error"] is None


def test_invalid_output_retries() -> None:
    # confidence outside [0.0, 1.0] should trigger a retry
    raw = (
        '{"type":"final_answer","content":"done",'
        '"confidence": 2.5}'
    )
    state = make_test_state(
        _raw_model_output=raw,
        _validation_attempt=0,
    )
    result = validation_node(state)
    assert result["_validation_attempt"] == 1
    assert "confidence must be" in result["_validation_error"]


def test_ceiling_breach_escalates() -> None:
    raw = (
        '{"type":"final_answer","content":"done",'
        '"confidence": 2.5}'
    )
    state = make_test_state(
        _raw_model_output=raw,
        _validation_attempt=3,
        _policy={"max_cycles": 10,
                 "token_ceiling": 10000,
                 "wall_clock_timeout_s": 60,
                 "cost_cap_usd": 1.0,
                 "validation_max_retries": 3},
    )
    # First, run the node to produce the error and increment attempt
    node_result = validation_node(state)
    merged = {**state, **node_result}
    # Then route and build escalation update
    assert route_after_validation(merged) == "escalate"
\end{lstlisting}

These tests exercise the core invariants: valid output clears the error fields,
invalid output increments the attempt counter, and exceeding the retry ceiling
routes to escalation. Chapter~28 will add graph-level tests that assert on the
interaction between Reason, Validation, and Action over multiple cycles.
